\chapter{绪论}

\section{研究背景和研究意义}

在互联网的时代，数据库系统是信息基础设施的核心组件，几乎支撑了所有关键业务系统的运行。
从在线交易、社交网络到云计算与人工智能应用，海量数据的高效存储、实时访问与可靠管理，直接决定了上层业务的性能、可用性与安全性。

随着互联网业务规模的持续扩大，分布式和云原生架构快速普及，传统的集中式数据库在扩展性、可用性以及运维成本方面逐渐暴露出难以克服的瓶颈，
对于满足大规模在线服务与大数据业务的需求，显得十分乏力。
分布式数据库在这样的背景下，应运而生。
通过在多节点上对数据进行分片与复制，实现了水平扩展与高可用，已成为支撑大规模在线服务与大数据业务的关键技术路径。


\section{主流数据库的不足}
面对分布式系统和云原生架构的挑战，主流数据库产品在扩展性、可用性和运维成本方面存在一些不足。
\paragraph{传统关系型数据库}
传统关系型数据库（RDS，如 MySQL、PostgreSQL）是企业级数据存储的基石，
以其严格的 ACID 特性、强大的 SQL 查询能力和成熟的生态系统，在金融、电商等核心业务场景中占据统治地位。
但是，传统 RDS 的架构设计初衷是基于“单机”或“主从”模型，这在分布式层面带来了一些本质性的架构缺憾。

RDS 主要依赖垂直扩展（Scale-up）来提升性能，虽然通过读写分离可以扩展读能力，但在写性能扩展和海量数据存储方面存在显著瓶颈。
为了应对海量数据，RDS 通常采用分库分表（Sharding）策略，但这破坏了关系型数据库的核心优势：
\begin{itemize}
\item 复杂查询受限：一旦数据被切分到不同的物理节点，跨库的 JOIN 操作、GROUP BY 聚合查询变得极其困难甚至不可用，往往需要在应用层或中间件中进行低效的数据拼装。
\item 全局唯一性挑战：在单机环境下简单的自增主键（Auto-increment ID），在分布式环境下失效，需要引入分布式 ID 生成器（如雪花算法），增加了系统的复杂性。
\end{itemize}

同时，RDS 在分布式事务的处理上显得力不从心，尤其是对网络不稳定环境下的支持：
传统 RDS 处理跨库事务通常依赖 XA 协议（两阶段提交，2PC）。
这是一个阻塞性协议，协调者需要等待所有参与者响应。
在网络不稳定的分布式环境下，XA 事务会导致严重的锁竞争和长尾延迟，系统吞吐量会随节点数增加而急剧下降。

\paragraph{Redis}
Redis 是高性能 KV 存储和共享缓存事实上的标准，以其亚毫秒级的响应速度著称，
凭借其高性能、高可用性及易用性，已成为互联网产品中实现共享缓存的首选中间件解决方案。。
然而，作为一款优先追求可用性（Availability）和分区容错性（Partition Tolerance）的系统，
Redis集群（Cluster）在分布式强一致性和全局事务层面存在显著的架构取舍。

Redis 集群采用无中心化的分片架构，虽然实现了水平扩展，但在处理复杂分布式逻辑时面临“分片隔离”的难题：
\begin{itemize}
    \item \textbf{原子性操作无法跨越分片边界：\\}
    Redis 的事务机制（包括 Lua 脚本和 MULTI/EXEC）严格局限于单分片（Single Slot）。
    无法在同一个“事务”中对不同分片上的键进行操作。
    \item \textbf{聚合与多键操作受限：\\}
    在单机模式下极其高效的集合运算（如 SUNION、SDIFF）或批量操作（MGET），
    在集群模式下如果涉及多个分片的 Key，通常无法直接执行。客户端必须自行实现低效的“分散-查询-合并”逻辑，失去了 Redis 原本“计算向数据移动”的高效特性。
\end{itemize}

\paragraph{其他新型数据库}
由于分布式和云原生架构的普及，涌现出了一批新型数据库产品，它们更注重云原生、存算分离以及对分布式事务的支持程度。

首先是支持强一致性与分布式ACID事务的NewSQL数据库。它们最显著的特点是，即使数据分布在不同的物理节点上，依然支持跨行、跨表的复杂事务的原子性操作。
但是它们一般只提供快照隔离（Snapshot Isolation），甚至于可串行化（Serializable）隔离级别。其中的典型代表有FoundationDB、TiKV等。
这类数据库通常是有中心化的，依赖中心节点进行事务的编排。
这些数据库的对事务的处理能力局限在很高的隔离级别上，在高并发、高吞吐量以及对一致性要求较低的场景下，性能表现不佳。

另外还有一类数据库强调高性能与可调一致性。这类数据库通常源自 Dynamo 论文\cite{Dynamo}的理念，
优先保证高可用性（Availability）和分区容错性（Partition tolerance），即 AP 系统。
但现代版本中加入了“轻量级事务”来弥补一致性的不足。它们一般支持轻量级的事务，只能保证最终一致性。其中的典型代表是Cassandra，DynamoDB等。
它们在高可用性和分区容错性方面表现出色，但底层架构上的限制，导致其在实现事务处理方面存在一些不足。
\\

综上所述，传统的数据库产品在分布式系统和云原生架构下存在一些不足，同时新型的数据库一般会在高可用性和高一致性方面选择其一作为设计目标。
构建一套兼顾高可用性、高吞吐量和可调一致性的分布式数据库系统，让用户根据实际需求灵活选择隔离级别，
从而取得性能和一致性之间的平衡是一个具有重要意义的研究方向，也是本文的工作目标。

\section{本文主要工作}
本文首先利用\TLAPlus 对分布式KV数据库的操作进行建模，分析不同隔离级别下事务处理的开始、过程中与结束后系统需要满足的安全属性。
然后，在规约建模的基础上，分析并设计了一套分布式KV数据库的实现方案，在利用\TLAPlus 的模型检测工具进行验证，确保算法的正确性。
之后，基于设计，利用SpringBoot-Dubbo-Nacos框架，使用docker技术模拟网络分区，给出了一种原型实现，
并在原型设计的基础上，进行了性能测试，验证了算法相对于强一致性模型在性能上的优势，相对于高可用性模型对一致性的提升。
实验证明，所设计的分布式KV数据库算法在性能和一致性之间取得了一定的平衡，满足了实际应用需求。

本文的主要工作内容如下：
\begin{itemize}
\item 利用\TLAPlus 对分布式KV数据库的操作进行建模，分析不同隔离级别下事务处理的开始、过程中与结束后系统需要满足的安全属性。
\item 在规约建模的基础上，分析并设计了一套分布式KV数据库的实现方案，在利用\TLAPlus 的模型检测工具进行验证，确保算法的正确性。
\item 基于设计，利用SpringBoot-Dubbo-Nacos框架，使用docker技术模拟网络分区，给出了一种原型实现。
\item 在原型实现的基础上进行性能测试，验证了算法相对于强一致性模型在性能上的优势，相对于高可用性模型对一致性的提升。
\end{itemize}

本文的组织结构如下：
%TODO
